\documentclass[11pt,a4paper,oneside]{article}

\usepackage[a4paper,left=3cm,right=2cm,top=2.5cm,bottom=2.5cm]{geometry}
\usepackage{tabto}
\usepackage[french]{babel}
\usepackage[T1]{fontenc}
\usepackage[utf8]{inputenc}

\begin{document}

\begin{titlepage}

  \vspace*{5.5cm}

  \begin{center}


      \noindent {\Huge Dossier de candidature à un poste d'Attaché Temporaire d'Enseignement et de Recherche}

      \vspace{1.5cm}

     {\huge Section 27}

      \vspace{2.5cm}
    \end{center}

    \begin{center}
        \noindent Maxime Gaide

    \end{center}

    \vspace{10cm}

    \noindent Doctorant au Laboratoire d'Informatique et d'Automatique pour les Systèmes

    \noindent ISAE-ENSMA / Université de Poitiers

    \noindent 2024


\end{titlepage}



\newpage

\section{Informations générales}

\noindent Nom Complet : Maxime Gaide  \\
Âge : 27 ans \\
Situation actuelle : Doctorant en troisième année \\
Laboratoire : Laboratoire d'Informatique et d'Automatique pour les Systèmes (LIAS)\\
\tabto{2.3cm}  ISAE-ENSMA / Université de Poitiers \\
Thématique de recherche : Informatique graphique analyse et reconstruction de formes \\
Section CNU : 27 \\
Adresse postale : 4 IMP DE LA HAUTE PAYRE, 86130 Jaunay-Marigny, France\\
Téléphone : +33786269524\\
Courriel : maxime.gaide@protonmail.ch \\
Langues :
\begin{itemize}
  \item Français : langue maternelle
  \item Anglais : professionnel scientifique écrit et oral
\end{itemize}

\section{Formation initiale et expériences}

\subsection{Formation Initiale}

\subsubsection{Licence Informatique - Aix-Marseille Université (2014-2018)}

La première année de licence permet d'acquérir les fondamentaux de
l'informatique avec des bases en mathématiques. %
Les modules d'introduction à la programmation sont enseignés en langage C. %
Les deux années suivantes forment, entre autres, à l'algorithmique, la théorie
des graphes, la programmation orienté objet (Java) ou encore à la compréhension
des systèmes informatiques. %
Cette formation offre également la possibilité de personnaliser ses compétences
via des modules optionnels (programmation fonctionnelle, connaissance de la
recherche en informatique). %

\subsubsection{Master 1 Informatique - Université d'Uppsala, Suède, (2018-2019)}

Cette année a été réalisée en Suède via une mobilité Erasmus+ comprenant un test
d'anglais au début ainsi qu'à la fin. %
Pour cette année, j'ai suivi des modules d'introduction à la programmation
parallèle, interaction Homme-Machine, base de données, gestion de projet (avec
offre d’intégration de fonctionnalité auprès d'autres groupes d'étudiants) et une
introduction à l'informatique graphique. %

\subsubsection{Master 2 Géométrie et Informatique Graphique - Aix-Marseille Université (2019-2020)}

La deuxième année se spécialise dans la modélisation 3D, l'informatique
graphique et la géométrie appliquée. %
Le semestre 3 forme aux fondamentaux de la géométrie discrète, de l'animation, à
l'analyse et au traitement de la géométrie. %
Ce parcours de Master étant indifférent vis-à-vis d'un parcours recherche ou
professionnel propose également un module autour réalisation d'un projet typé
industriel avec un cahier des charges et sa documentation. %
Le semestre 4 est quant à lui partagé entre un projet de recherche en groupe au sein du Laboratoire d'Informatique et Système (LIS) de Marseille et d'un stage (recherche ou industriel).

\subsubsection{Doctorat en Informatique - Rejeu et modélisation spatio-temporelle à base de règles d'éléments archéologiques  (2021-en cours)}
Directeur de thèse : Stéphane Jean\\
Encadrants : David Marcheix, Xavier Skapin\\

\subsection{Stages et expériences}

\subsubsection{Stage recherche - Détection de cratères dans les maillages surfaciques de corps célestes - Aix-Marseille Université (2019, 2 mois)}


\subsubsection{Stage recherche - Compression de pipeline 3D pour le streaming d'application - Nothing2Install, Marseille, 2020, 6 mois)}

\subsubsection{Surveillance d'examens, Toulon, (2021)}

\section{Enseignements}

\subsection{Algorithmique et Structures Numériques; 1ère année ENSMA; 27h TP, 45h Projet, 3h TD}

\subsection{Base de la Conception Logicielle; 1ère année ENSMA; 42h TP}
\subsection{Utilisation et Exploitation des Données; 1ère année ENSMA; 6h TP}
\subsection{Apprentissage Automatique appliqué; 2ème année ENSMA; 3h TP}
\subsection{Récapitulatif}
Le tableau récapitulatif suivant, synthétise les données précédentes :

\begin{center}
\begin{tabular}{|l|l|l|l|l|l|}
\hline
Établissement & Matière Enseignée & Niveau & \multicolumn{3}{l|}{Nombre d'heures} \\
\hline
 & & & TD & TP & Projet \\
\hline
ENSMA & Algorithmique et Systèmes Numériques & A1 ingénieur & 3h45 & 27h & 45h \\
\hline
ENSMA & Base de la Conception Logicielle & A1 ingénieur & & 42h & \\
\hline
ENSMA & Utilisation et Exploitation des Données & A1 ingénieur & & 6h & \\
\hline
ENSMA & Apprentissage automatique appliqué & A2 ingénieur & & 3h & \\
\hline
\multicolumn{3}{|l|}{Total:} & \multicolumn{3}{l|}{126h45} \\
\hline
\multicolumn{3}{|l|}{Total heures équivalent TD:} & \multicolumn{3}{l|}{126h45} \\
\hline
\end{tabular}
\end{center}

\section{Recherche}

\subsection{Analyses syntaxiques de règles de transformations de graphe}
\subsection{Réévaluation de modèle à partir de règles de transformation de graphe}
\subsection{ Exploitation des scripts de règles pour la réévaluation de modèles}
\end{document}
